% ══════════════════════════════════════════════════════════════════════════════
% Lecture 9 : Programmation Dynamique – Coupe de tiges (reconstruction)
%             et Multiplication de chaînes de matrices
% ══════════════════════════════════════════════════════════════════════════════

% ─────────────────────────────────────────────────────────────────────────────
\subsection{Rappel : éléments clés d'un algorithme DP}

La programmation dynamique repose sur deux propriétés fondamentales :

\begin{enumerate}
  \item \textbf{Sous-structure optimale} : une solution optimale s'obtient en faisant
        un \emph{choix}, ce qui laisse un ou plusieurs sous-problèmes à résoudre,
        et la solution optimale globale résout ces sous-problèmes de façon optimale.
  \item \textbf{Sous-problèmes répétés} : un algorithme récursif naïf recalcule les mêmes
        sous-problèmes de nombreuses fois.
        \begin{itemize}
          \item \textbf{Top-down avec mémoïsation} : résoudre récursivement en
                mémorisant chaque résultat dans une table.
          \item \textbf{Bottom-up} : trier les sous-problèmes et résoudre d'abord
                les plus petits ; ainsi, quand on résout un sous-problème, les
                sous-problèmes dont il dépend sont déjà calculés.
        \end{itemize}
\end{enumerate}

% ─────────────────────────────────────────────────────────────────────────────
\subsection{Coupe de tiges – Reconstruction de la solution optimale}

En Lecture~8, l'algorithme DP renvoyait uniquement le revenu optimal $r(n)$.
Il est souvent nécessaire de récupérer \emph{comment} obtenir ce revenu, c'est-à-dire
les longueurs des morceaux à couper.

\subsubsection{Algorithme étendu avec table des décisions}

L'idée est de maintenir une table auxiliaire $s[0\ldots n]$ : $s[j]$ stocke la longueur
du premier morceau coupé (la \emph{décision}) correspondant au revenu optimal $r[j]$.

\begin{algorithm}[H]
\caption{Extended-Bottom-Up-Cut-Rod($p$, $n$)}
\begin{algorithmic}[1]
\State Créer les tableaux $r[0\ldots n]$ et $s[0\ldots n]$
\State $r[0] \leftarrow 0$
\For{$j = 1$ \textbf{to} $n$}
    \State $q \leftarrow -\infty$
    \For{$i = 1$ \textbf{to} $j$}
        \If{$q < p[i] + r[j - i]$}
            \State $q \leftarrow p[i] + r[j - i]$  \AlgComment{meilleur revenu}
            \State $s[j] \leftarrow i$              \AlgComment{stocker le choix optimal}
        \EndIf
    \EndFor
    \State $r[j] \leftarrow q$
\EndFor
\State \Return $r$ et $s$
\end{algorithmic}
\end{algorithm}

La complexité reste $\Theta(n^2)$. La seule différence avec l'algorithme de base
est le maintien de $s[j]$ qui enregistre l'indice $i$ réalisant le maximum.

\subsubsection{Exemple : tables $r$ et $s$ pour $n = 8$}

En utilisant la table de prix $p_i$ : 1, 5, 8, 9, 10, 17, 17, 20, 24, 30 :

\begin{center}
\renewcommand{\arraystretch}{1.3}
\begin{tabular}{ccccccccccc}
\toprule
$i$    & 0 & 1 & 2 & 3 & 4  & 5  & 6  & 7  & 8  \\
\midrule
$r[i]$ & 0 & 1 & 5 & 8 & 10 & 13 & 17 & 18 & 22 \\
$s[i]$ & 0 & 1 & 2 & 3 & 2  & 2  & 6  & 1  & 2  \\
\bottomrule
\end{tabular}
\end{center}

\subsubsection{Extraction de la solution}

La table $s$ stocke les choix menant à une solution optimale.
On remonte la solution en traçant-back depuis $n$ :

\begin{algorithm}[H]
\caption{Print-Cut-Rod-Solution($p$, $n$)}
\begin{algorithmic}[1]
\State $(r, s) \leftarrow \textsc{Extended-Bottom-Up-Cut-Rod}(p, n)$
\While{$n > 0$}
    \State \textbf{afficher} $s[n]$   \AlgComment{taille du prochain morceau}
    \State $n \leftarrow n - s[n]$
\EndWhile
\end{algorithmic}
\end{algorithm}

\textbf{Trace pour $n = 8$ :}

\begin{center}
\begin{tikzpicture}[font=\small, >=Stealth]
  \foreach \n/\out/\x in {8/{}/0, 6/{2}/2.5, 0/{6}/5}{
    \node[draw, fill=blue!8, rounded corners, minimum width=1.8cm,
          minimum height=0.6cm, align=center] at (\x, 0) {$n = \n$};
    \ifx\out\empty\else
      \node[draw=none, fill=none, keygreen, font=\bfseries\small] at (\x-1.25, 0)
            {sortie : \out};
    \fi
  }
  \draw[->, thick, gray] (1.0, 0) -- (1.5, 0);
  \draw[->, thick, gray] (3.5, 0) -- (4.0, 0);
  \node[below=8pt, keyred, font=\small] at (2.5, -0.3)
        {Découpe : \textbf{2 + 6}, revenu $= 5 + 17 = 22$};
\end{tikzpicture}
\end{center}

% ─────────────────────────────────────────────────────────────────────────────
\subsection{Résumé : coupe de tiges}

\begin{tcolorbox}[colback=lightblue, colframe=keyblue, fonttitle=\bfseries,
                  title={Récapitulatif DP – Coupe de tiges}]
\begin{itemize}[noitemsep]
  \item \textbf{Choix} : où effectuer la coupe la plus à gauche.
  \item \textbf{Sous-structure optimale} : le reste doit être découpé de façon optimale.
  \item \textbf{Récurrence} :
    \[
      r(n) =
      \begin{cases}
        0 & \text{si } n = 0, \\
        \displaystyle\max_{1 \le i \le n} \bigl\{p_i + r(n-i)\bigr\} & \text{si } n \ge 1.
      \end{cases}
    \]
  \item \textbf{Complexité} : $\Theta(n^2)$ (top-down ou bottom-up).
  \item \textbf{Reconstruction} : table auxiliaire $s$ + trace-back en $O(n)$.
\end{itemize}
\end{tcolorbox}

% ─────────────────────────────────────────────────────────────────────────────
\subsection{Quand peut-on utiliser la programmation dynamique ?}

\begin{enumerate}
  \item \textbf{Sous-structure optimale.}
    \begin{itemize}
      \item Une solution optimale peut être construite en combinant des solutions
            optimales de sous-problèmes.
      \item Cela implique que la valeur optimale peut être exprimée par une
            \emph{formule récursive}.
    \end{itemize}
  \item \textbf{Sous-problèmes répétés.}
    \begin{itemize}
      \item Un algorithme récursif naïf peut revisiter le même sous-problème
            de nombreuses fois.
    \end{itemize}
\end{enumerate}

\begin{remark}{Problème du rendu de monnaie}{change9}
La DP s'applique directement à ce problème classique. Avec $n$ dénominations
$0 < w_1 < \cdots < w_n$ et un montant cible $W$, on minimise le nombre de pièces :
\[
  \min\left\{\sum_{j=1}^{n} x_j \;:\; \sum_{j=1}^{n} w_j x_j = W,\; x_j \in \mathbb{Z}_{\ge 0}\right\}.
\]
\textbf{Exemple :} $w_1=1$, $w_2=2$, $w_3=5$, $W=8$.
Optimal : $x_1 = x_2 = x_3 = 1 \Rightarrow$ 3~pièces.
\end{remark}

% ══════════════════════════════════════════════════════════════════════════════
\section*{Multiplication de chaînes de matrices}
\addcontentsline{toc}{subsection}{Multiplication de chaînes de matrices}
% ══════════════════════════════════════════════════════════════════════════════

% ─────────────────────────────────────────────────────────────────────────────
\subsection{Coût de la multiplication de deux matrices}

Multiplier une matrice $A_{p \times q}$ par une matrice $B_{q \times r}$ produit une
matrice $C_{p \times r}$.

\begin{center}
\begin{tikzpicture}[font=\small, >=Stealth]
  % Matrix A
  \draw[fill=blue!10, draw=keyblue] (0,0) rectangle (1.4, 2.0);
  \node at (0.7, 2.3) {$q$};
  \node[left] at (-0.05, 1.0) {$p$};
  \node at (0.7, 1.0) {$A_{p,q}$};
  \node[font=\Large] at (1.9, 1.0) {$\times$};
  % Matrix B
  \draw[fill=orange!15, draw=orange!70!black] (2.4,0.3) rectangle (4.4, 2.0);
  \node at (3.4, 2.3) {$r$};
  \node[left] at (2.35, 1.15) {$q$};
  \node at (3.4, 1.15) {$B_{q,r}$};
  \node[font=\Large] at (4.9, 1.0) {$=$};
  % Matrix C
  \draw[fill=green!12, draw=keygreen] (5.4, 0) rectangle (7.4, 2.0);
  \node at (6.4, 2.3) {$r$};
  \node[left] at (5.35, 1.0) {$p$};
  \node at (6.4, 1.0) {$C_{p,r}$};
  % Annotation
  \node[below=4pt, keyred, font=\bfseries\small] at (3.5, -0.2)
        {Coût : $p \cdot q \cdot r$ multiplications scalaires};
\end{tikzpicture}
\end{center}

Chaque entrée de $C$ nécessite $q$ multiplications scalaires.
La matrice $C$ contient $p \cdot r$ entrées, donc le coût total est
$\mathbf{p \cdot q \cdot r}$ \textbf{multiplications scalaires}.

\begin{tcolorbox}[colback=lightyellow, colframe=orange!70!black, fonttitle=\bfseries,
                  title={Pourquoi l'ordre d'association importe-t-il ?}]
Pour $A$ ($20 \times 1$), $B$ ($1 \times 10$), $C$ ($10 \times 100$) :
\begin{itemize}[noitemsep]
  \item $(A \times (B \times C))$ : $1\cdot10\cdot100 + 20\cdot1\cdot100 = 3\,000$
  \item $((A \times B) \times C)$ : $20\cdot1\cdot10 + 20\cdot10\cdot100 = 20\,200$
\end{itemize}
La parenthésisation peut \textbf{modifier radicalement} le nombre d'opérations.
\end{tcolorbox}

% ─────────────────────────────────────────────────────────────────────────────
\subsection{Définition du problème}

\begin{definition}{Multiplication de chaînes de matrices}{mcm}
\textbf{Entrée :} une chaîne $\langle A_1, A_2, \ldots, A_n \rangle$ de $n$ matrices,
où la matrice $A_i$ a la dimension $p_{i-1} \times p_i$ pour $i = 1, \ldots, n$.

\textbf{Sortie :} une parenthésisation complète du produit $A_1 A_2 \cdots A_n$
minimisant le nombre de multiplications scalaires.
\end{definition}

\textbf{Remarques :}
\begin{itemize}
  \item On ne demande \emph{pas} de calculer le produit, seulement de trouver la
        meilleure parenthésisation.
  \item La parenthésisation peut affecter considérablement le nombre de
        multiplications.
\end{itemize}

\textbf{Exemple :} chaîne $A_1 A_2 A_3$ avec dimensions $50\times5$, $5\times100$,
$100\times10$ :

\begin{center}
\renewcommand{\arraystretch}{1.3}
\begin{tabular}{lrl}
\toprule
\textbf{Parenthésisation} & \textbf{Coût} & \textbf{Calcul} \\
\midrule
$(A_1 A_2) A_3$ & $75\,000$ & $50\cdot5\cdot100 + 50\cdot100\cdot10$ \\
$A_1 (A_2 A_3)$ & $7\,500$  & $5\cdot100\cdot10 + 50\cdot5\cdot10$ \\
\bottomrule
\end{tabular}
\end{center}

% ─────────────────────────────────────────────────────────────────────────────
\subsection{Sous-structure optimale}

\begin{theorem}{Sous-structure optimale – Multiplication de matrices}{mcmopt}
Supposons que la parenthésisation optimale de $A_1 A_2 \cdots A_n$ effectue la
dernière multiplication comme $(A_1 \cdots A_k)(A_{k+1} \cdots A_n)$. Alors les
parenthésisations de $A_1 \cdots A_k$ et $A_{k+1} \cdots A_n$ dans cette solution
optimale sont elles-mêmes optimales.
\end{theorem}

\begin{proof}
Soit $M(P)$ le nombre de multiplications scalaires d'une parenthésisation $P$.
Notons $((O_L) \cdot (O_R))$ la solution optimale où $O_L$ et $O_R$ parenthésisent
respectivement $A_1\cdots A_k$ et $A_{k+1}\cdots A_n$. Soient $P_L$ et $P_R$ des
parenthésisations optimales de ces deux sous-chaînes. On a :
\begin{align*}
  M((O_L) \cdot (O_R))
    &= p_0 \cdot p_k \cdot p_n + M(O_L) + M(O_R) \\
    &\ge p_0 \cdot p_k \cdot p_n + M(P_L) + M(P_R)
     = M((P_L) \cdot (P_R)).
\end{align*}
Puisque $((O_L)\cdot(O_R))$ est optimal, on a l'égalité, donc
$M(P_L) = M(O_L)$ et $M(P_R) = M(O_R)$ : $P_L$ et $P_R$ sont bien optimales.
\end{proof}

% ─────────────────────────────────────────────────────────────────────────────
\subsection{Formule récursive}

Soit $m[i, j]$ le nombre minimal de multiplications scalaires pour calculer
$A_i A_{i+1} \cdots A_j$.

\[
  m[i, j] =
  \begin{cases}
    0 & \text{si } i = j, \\
    \displaystyle\min_{i \le k < j}
      \bigl\{ m[i,k] + m[k+1,j] + p_{i-1}\, p_k\, p_j \bigr\}
    & \text{si } i < j.
  \end{cases}
\]

\begin{remark}{Ordre de remplissage}{ordrem}
$m[i, j]$ ne dépend que de sous-problèmes avec un intervalle $j - i$ \emph{plus petit}.
Un algorithme bottom-up résout donc les sous-problèmes par ordre croissant de $j - i$
(longueur de chaîne $\ell = j - i + 1$).
\end{remark}

% ─────────────────────────────────────────────────────────────────────────────
\subsection{Algorithme bottom-up}

\begin{algorithm}[H]
\caption{Matrix-Chain-Order($p$)}
\begin{algorithmic}[1]
\State $n \leftarrow p.\mathit{length} - 1$
\State Créer les tables $m[1\ldots n,\, 1\ldots n]$ et $s[1\ldots n,\, 1\ldots n]$
\For{$i = 1$ \textbf{to} $n$}
    \State $m[i, i] \leftarrow 0$
\EndFor
\For{$\ell = 2$ \textbf{to} $n$}                  \AlgComment{$\ell$ = longueur de la chaîne}
    \For{$i = 1$ \textbf{to} $n - \ell + 1$}
        \State $j \leftarrow i + \ell - 1$
        \State $m[i, j] \leftarrow \infty$
        \For{$k = i$ \textbf{to} $j - 1$}
            \State $q \leftarrow m[i,k] + m[k+1,j] + p_{i-1}\, p_k\, p_j$
            \If{$q < m[i,j]$}
                \State $m[i,j] \leftarrow q$
                \State $s[i,j] \leftarrow k$   \AlgComment{$s$ mémorise le choix optimal}
            \EndIf
        \EndFor
    \EndFor
\EndFor
\State \Return $m$ et $s$
\end{algorithmic}
\end{algorithm}

\begin{complexity}{Matrix-Chain-Order}{mcmcplx}
\begin{itemize}[noitemsep]
  \item La table $m$ contient $O(n^2)$ entrées.
  \item Chaque entrée $m[i,j]$ requiert $O(n)$ pour trouver le meilleur $k$.
  \item \textbf{Complexité totale : $\Theta(n^3)$ en temps, $\Theta(n^2)$ en espace.}
\end{itemize}
\end{complexity}

% ─────────────────────────────────────────────────────────────────────────────
\subsection{Exemple complet}

Chaîne de $n = 6$ matrices avec dimensions :

\begin{center}
\renewcommand{\arraystretch}{1.2}
\begin{tabular}{lcccccc}
\toprule
Matrice     & $A_1$          & $A_2$          & $A_3$         & $A_4$        & $A_5$          & $A_6$ \\
\midrule
Dimensions & $30\times35$ & $35\times15$ & $15\times5$ & $5\times10$ & $10\times20$ & $20\times25$ \\
\bottomrule
\end{tabular}
\end{center}

\noindent Soit $p = \langle 30, 35, 15, 5, 10, 20, 25 \rangle$.

Les tables $m$ (coûts optimaux) et $s$ (choix optimaux) calculées par
\textsc{Matrix-Chain-Order} sont :

\begin{center}
\renewcommand{\arraystretch}{1.4}
\begin{tabular}{c|cccccc}
\multicolumn{7}{c}{\textbf{Table $m[i,j]$}} \\
\hline
$j \backslash i$ & 1 & 2 & 3 & 4 & 5 & 6 \\
\hline
1 & 0 & & & & & \\
2 & 15\,750 & 0 & & & & \\
3 & 7\,875 & 2\,625 & 0 & & & \\
4 & 9\,375 & 4\,375 & 750 & 0 & & \\
5 & 11\,875 & 7\,125 & 2\,500 & 1\,000 & 0 & \\
6 & 15\,125 & 10\,500 & 5\,375 & 3\,500 & 5\,000 & 0 \\
\hline
\end{tabular}
\qquad
\begin{tabular}{c|cccccc}
\multicolumn{7}{c}{\textbf{Table $s[i,j]$}} \\
\hline
$j \backslash i$ & 1 & 2 & 3 & 4 & 5 & 6 \\
\hline
2 & 1 & & & & & \\
3 & 1 & 2 & & & & \\
4 & 3 & 3 & 3 & & & \\
5 & 3 & 3 & 3 & 4 & & \\
6 & 3 & 3 & 3 & 5 & 5 & \\
\hline
\end{tabular}
\end{center}

La solution optimale est $m[1,6] = 15\,125$ multiplications scalaires.
En remontant depuis $s[1,6] = 3$, on obtient la parenthésisation :
\[
  \bigl(A_1 \,(A_2\,A_3)\bigr)\bigl((A_4\,A_5)\,A_6\bigr)
\]

% ─────────────────────────────────────────────────────────────────────────────
\subsection{Reconstruction de la parenthésisation optimale}

La table $s$ permet de reconstruire récursivement la parenthésisation :

\begin{algorithm}[H]
\caption{Print-Optimal-Parens($s$, $i$, $j$)}
\begin{algorithmic}[1]
\If{$i = j$}
    \State \textbf{afficher} ``$A_i$''
\Else
    \State \textbf{afficher} ``(''
    \State \Call{Print-Optimal-Parens}{$s$, $i$, $s[i,j]$}
    \State \Call{Print-Optimal-Parens}{$s$, $s[i,j]+1$, $j$}
    \State \textbf{afficher} ``)''
\EndIf
\end{algorithmic}
\end{algorithm}

\textbf{Trace pour l'exemple ($i=1$, $j=6$) :}

\begin{center}
\begin{tikzpicture}[font=\small, >=Stealth,
    node/.style={draw, rounded corners, fill=blue!8,
                 minimum width=2.6cm, minimum height=0.55cm,
                 inner sep=4pt}]
  \node[node] (r) at (0, 0) {$s[1,6]=3$};
  \node[node, fill=green!12] (l) at (-3, -1.1) {$s[1,3]=1 \Rightarrow A_1(A_2 A_3)$};
  \node[node, fill=orange!15] (r2) at (3, -1.1) {$s[4,6]=5 \Rightarrow (A_4 A_5)A_6$};
  \draw[->] (r) -- (l);
  \draw[->] (r) -- (r2);
  \node[below=16pt, keyred] at (0, -1.5)
        {Résultat : $\bigl(A_1\,(A_2\,A_3)\bigr)\bigl((A_4\,A_5)\,A_6\bigr)$};
\end{tikzpicture}
\end{center}

% ─────────────────────────────────────────────────────────────────────────────
\subsection{Résumé : multiplication de chaînes de matrices}

\begin{tcolorbox}[colback=lightblue, colframe=keyblue, fonttitle=\bfseries,
                  title={Récapitulatif DP – Multiplication de chaînes de matrices}]
\begin{itemize}[noitemsep]
  \item \textbf{Choix} : où placer la parenthèse la plus externe
        $(A_1 \cdots A_k)(A_{k+1} \cdots A_n)$.
  \item \textbf{Sous-structure optimale} : les deux sous-chaînes doivent être
        parenthésées de façon optimale (prouvé par l'argument de cut-and-paste).
  \item \textbf{Récurrence} :
    \[
      m[i,j] =
      \begin{cases}
        0 & \text{si } i = j, \\
        \displaystyle\min_{i \le k < j}
          \bigl\{m[i,k] + m[k+1,j] + p_{i-1}\,p_k\,p_j\bigr\}
        & \text{si } i < j.
      \end{cases}
    \]
  \item \textbf{Complexité} : $\Theta(n^3)$ en temps, $\Theta(n^2)$ en espace.
  \item \textbf{Reconstruction} : table auxiliaire $s$ + \textsc{Print-Optimal-Parens}.
\end{itemize}
\end{tcolorbox}

% ─────────────────────────────────────────────────────────────────────────────
\subsection{Recette générale pour la DP (consolidation)}

\begin{center}
\renewcommand{\arraystretch}{1.45}
\begin{tabular}{llcl}
\toprule
\textbf{Problème} & \textbf{Choix} & \textbf{Complexité} & \textbf{Taille de la table} \\
\midrule
Coupe de tiges     & Position de la coupe gauche & $\Theta(n^2)$   & $n+1$ entrées \\
Rendu de monnaie   & Dénomination utilisée       & $O(nW)$         & $W+1$ entrées \\
Chaînes de matrices & Position de la parenthèse  & $\Theta(n^3)$   & $n^2$ entrées \\
\bottomrule
\end{tabular}
\end{center}

Les trois étapes systématiques restent :
\begin{enumerate}
  \item Identifier les choix et la sous-structure optimale.
  \item Écrire la récurrence pour la valeur optimale.
  \item Résoudre efficacement (top-down avec mémoïsation ou bottom-up).
\end{enumerate}